\section{Stack Frame, Konvensi Pemanggilan, dan Rekursi}
\label{sec:stack-frame}

Tabel~\ref{tab:trace-stack} memperlihatkan bahwa prosedur selalu
mengerjakan empat hal: membangun bingkai, mengakses parameter,
menghasilkan nilai balik, lalu membongkar bingkai. Pola bakunya
disebut \term{stack frame} standar \cite{bartlett2004}.

\subsection{\eng{Prologue} dan \eng{Epilogue}}

Bagian pembuka prosedur, atau \term{prologue}, terdiri atas dua
instruksi; bagian penutupnya, atau \term{epilogue}, juga dua
instruksi.

\begin{lstlisting}[
    language=MASM,
    caption={Kerangka baku \eng{prologue} dan \eng{epilogue} prosedur
    \eng{near}.},
    label={lst:prologue-epilogue}
]
HitungLuas PROC
    PUSH BP                 ; prologue : simpan BP pemanggil
    MOV  BP, SP             ; prologue : BP <- titik acuan bingkai

    ; --- badan prosedur ---
    ; parameter diakses sebagai [BP+4], [BP+6], ...

    POP  BP                 ; epilogue : pulihkan BP pemanggil
    RET                     ; epilogue : kembali ke pemanggil
HitungLuas ENDP
\end{lstlisting}

Instruksi \lstasm|PUSH BP| wajib ada karena \reg{BP} hanya ada satu
buah; tanpa menyimpannya, prosedur akan menimpa \reg{BP} milik
pemanggil sehingga akses parameter pemanggil menjadi salah.
Pasangan \lstasm|PUSH BP| dan \lstasm|POP BP| harus seimbang pada
setiap jalur keluar prosedur.

\subsection{Mengakses Parameter dan Variabel Lokal}

Setelah \lstasm|MOV BP, SP|, susunan bingkai selalu sama dengan
Gambar~\ref{fig:isi-stack-call}, sehingga jarak setiap isi dari
\reg{BP} bersifat tetap:

\begin{itemize}
    \item \code{[BP+0]} --- nilai \reg{BP} milik pemanggil;
    \item \code{[BP+2]} --- alamat kembali (\reg{IP});
    \item \code{[BP+4]} --- parameter pertama (yang didorong terakhir);
    \item \code{[BP+6]} --- parameter kedua;
    \item \code{[BP+8]} --- parameter ketiga, dan seterusnya.
\end{itemize}

Karena pemanggil mendorong parameter kedua lebih dahulu, parameter
pertama-lah yang berada pada \code{[BP+4]}. Bila urutan \asm{PUSH}
dibalik tanpa mengubah badan prosedur, kedua nilai akan tertukar
tanpa pesan kesalahan apa pun.

Variabel lokal disimpan di bawah \reg{BP}, misalnya pada
\code{[BP-2]}, dengan mengurangi \reg{SP} pada akhir
\eng{prologue} dan memulihkannya pada awal \eng{epilogue}:

\begin{lstlisting}[language=MASM, numbers=none]
    PUSH BP                 ; prologue
    MOV  BP, SP
    SUB  SP, 4              ; sediakan dua variabel lokal 16-bit

    MOV  WORD PTR [BP-2], 0 ; variabel lokal pertama
    MOV  WORD PTR [BP-4], 0 ; variabel lokal kedua

    MOV  SP, BP             ; epilogue : buang variabel lokal
    POP  BP
    RET
\end{lstlisting}

\subsection{Konvensi Pemanggilan}

\term{Konvensi pemanggilan} (\eng{calling convention}) adalah
kesepakatan antara pemanggil dan prosedur mengenai tiga hal: letak
parameter, letak nilai balik, dan siapa yang membersihkan stack.
Tanpa kesepakatan ini, prosedur yang ditulis satu orang tidak dapat
dipakai oleh program orang lain \cite{uva2024}. Tiga cara pengiriman
parameter yang lazim dipakai dirangkum pada Tabel~\ref{tab:konvensi}.

\begin{table}[H]
    \centering
    \caption{Perbandingan cara pengiriman parameter kepada prosedur.}
    \label{tab:konvensi}
    \small
    \begin{tabular}{@{}p{2.4cm}p{3.6cm}p{3.7cm}p{3.5cm}@{}}
        \toprule
        \textbf{Cara} & \textbf{Kelebihan} &
        \textbf{Kekurangan} & \textbf{Contoh} \\
        \midrule
        Melalui register & Tidak ada akses memori; paling cepat dan
        paling ringkas & Jumlah parameter terbatas oleh jumlah
        register; prosedur menjadi kurang umum &
        \lstasm|MOV AX, 5| lalu \lstasm|CALL Tulis| \\
        \addlinespace
        Melalui stack & Jumlah parameter tidak terbatas; pemanggil
        tidak perlu tahu register mana yang dipakai prosedur &
        Lebih lambat karena setiap parameter menimbulkan akses
        memori; rawan stack tidak seimbang &
        \lstasm|PUSH 5| lalu \lstasm|CALL HitungLuas| \\
        \addlinespace
        Melalui variabel global & Prosedur sangat sederhana; tidak
        ada instruksi tambahan sebelum pemanggilan & Prosedur tidak
        dapat dipakai ulang; sulit ditelusuri saat kesalahan terjadi
        & \lstasm|MOV lebar, 5| lalu \lstasm|CALL HitungLuas| \\
        \bottomrule
    \end{tabular}
\end{table}

Siapa yang membersihkan stack dari parameter? Bila prosedur yang
membersihkannya, \eng{epilogue} berakhir dengan \asm{RET} \code{n}
yang mengambil \reg{IP} lalu menambahkan \code{n} pada \reg{SP}
sekaligus. Bila pemanggil yang membersihkannya, prosedur memakai
\asm{RET} biasa dan pemanggil menambahkan \lstasm|ADD SP, 4|
setelah \asm{CALL}. Keduanya tidak boleh dipakai bersama-sama:
\asm{RET} \code{4} yang disusul \lstasm|ADD SP, 4| membuat \reg{SP}
naik \dec{8} byte padahal hanya \dec{4} byte yang perlu dibuang.

Pertanyaan ketiga menyangkut register. Pada 8086 tidak ada aturan
perangkat keras yang memaksa, sehingga kesepakatan dibuat di tingkat
program; pembagian yang lazim mengikuti Microsoft C untuk 8086
dirangkum pada Tabel~\ref{tab:register-simpan}.

\begin{table}[H]
    \centering
    \caption{Pembagian tanggung jawab penyimpanan register pada
    konvensi pemanggilan 8086.}
    \label{tab:register-simpan}
    \small
    \begin{tabular}{@{}p{3.3cm}p{2.6cm}p{3.5cm}p{3.8cm}@{}}
        \toprule
        \textbf{Golongan} & \textbf{Register} &
        \textbf{Tanggung jawab} & \textbf{Akibat bila dilanggar} \\
        \midrule
        Boleh dirusak (\eng{caller-saved}) & \reg{AX}, \reg{BX},
        \reg{CX}, \reg{DX}, \reg{ES} & Pemanggil menyimpan sendiri
        bila nilainya masih diperlukan setelah pemanggilan &
        Pemanggil memakai nilai yang sudah berubah \\
        \addlinespace
        Wajib dipulihkan (\eng{callee-saved}) & \reg{BP}, \reg{SI},
        \reg{DI}, \reg{DS}, \reg{SS} & Prosedur menyimpan pada
        \eng{prologue} dan memulihkan pada \eng{epilogue} &
        Pemanggil kehilangan nilai lama; akses data dapat salah
        alamat \\
        \addlinespace
        Selalu dipulihkan & \reg{SP} & Stack harus seimbang saat
        \asm{RET} dijalankan & \asm{RET} mengambil alamat yang salah;
        program berhenti tak terkendali \\
        \bottomrule
    \end{tabular}
\end{table}

\subsection{Nilai Balik Prosedur}

Nilai balik dikembalikan melalui register, bukan melalui stack:
hasil \dec{8} bit pada \reg{AL}, hasil \dec{16} bit pada \reg{AX},
dan hasil \dec{32} bit pada pasangan \reg{DX}:\reg{AX}. Pola
\reg{DX}:\reg{AX} muncul secara alami karena \asm{MUL} dan
\asm{IMUL} memang menyimpan hasil kalinya pada kedua register itu
(Bab~\ref{chap:aritmatika}). \reg{AX} termasuk register yang boleh
dirusak, sehingga prosedur tidak wajib memulihkannya.

\subsection{Rekursi pada Assembly}

\term{Rekursi} adalah pemanggilan prosedur oleh dirinya sendiri.
Pada bahasa tingkat tinggi rekursi tampak sebagai fitur bahasa; pada
assembly rekursi tampak jelas sebagai penumpukan bingkai stack.
Setiap pemanggilan menambah satu bingkai baru yang berisi parameter,
alamat kembali, dan \reg{BP} sendiri.

\begin{lstlisting}[
    language=MASM,
    caption={Prosedur faktorial rekursif (bagian inti).},
    label={lst:faktorial-inti}
]
Faktorial PROC
    PUSH BP                 ; prologue
    MOV  BP, SP

    MOV  AX, [BP+4]         ; AX <- n
    CMP  AX, 1              ; apakah n <= 1 ?
    JBE  basis              ; ya -> kasus basis

    DEC  AX                 ; AX <- n - 1
    PUSH AX                 ; dorong argumen (n-1)
    CALL Faktorial          ; rekursi : AX <- (n-1)!
    ADD  SP, 2              ; bersihkan argumen rekursif

    MUL  WORD PTR [BP+4]    ; DX:AX <- (n-1)! * n

    POP  BP                 ; epilogue
    RET

basis:
    MOV  AX, 1              ; 0! = 1! = 1
    POP  BP                 ; epilogue
    RET
Faktorial ENDP
\end{lstlisting}

Setiap bingkai rekursif memerlukan \dec{6} byte stack: \dec{2} byte
parameter, \dec{2} byte alamat kembali, dan \dec{2} byte \reg{BP}
lama. Untuk menghitung $5!$ diperlukan lima bingkai, yaitu sekitar
\dec{30} byte, sehingga kebutuhan stack tumbuh sebanding dengan
kedalaman rekursi atau setara $O(n)$.

Fibonacci rekursif $F(n) = F(n-1) + F(n-2)$ jauh lebih berat karena
setiap tingkat memanggil dirinya dua kali. Untuk $n = \dec{20}$ saja
jumlah pemanggilannya melebihi \dec{10.000}, sedangkan versi
iteratifnya cukup memakai satu perulangan dengan stack tetap. Karena
itu Fibonacci jauh lebih baik ditulis secara iteratif pada assembly.

\begin{peringatanbox}
Bahaya terbesar rekursi pada 8086 adalah \term{stack overflow}.
Prosesor 8086 tidak memiliki pengaman batas stack. Bila rekursi
terlalu dalam atau kasus basis tidak pernah tercapai, \reg{SP} akan
terus turun melewati batas bawah ruang \dir{.STACK}, lalu mulai
menimpa segmen data atau segmen kode. Kerusakannya terjadi tanpa
pesan kesalahan, dan program berhenti dengan perilaku yang sulit
ditelusuri. Pemeriksaan isi stack dengan \eng{debugger} dibahas pada
Bab~\ref{chap:debugging}.
\end{peringatanbox}

\subsection{Kesalahan Umum Pengelolaan Stack}

Empat kesalahan berikut paling sering dijumpai.

\begin{enumerate}
    \item \textbf{Stack tidak seimbang.} Jumlah \asm{PUSH} tidak
          sama dengan jumlah \asm{POP} pada setiap jalur eksekusi,
          sehingga \asm{RET} mengambil alamat yang salah.

    \item \textbf{\asm{RET} ke alamat yang salah.} Pemanggilan
          \eng{near} ditutup dengan \asm{RETF} atau sebaliknya,
          sehingga jumlah byte yang diambil dari stack tidak cocok.

    \item \textbf{Lupa menyimpan register yang dipakai.} Prosedur
          memakai \reg{SI} atau \reg{DI} tanpa menyimpannya,
          sehingga pemanggil kehilangan nilai lama register itu.

    \item \textbf{Kerusakan \eng{stack frame}.} \reg{BP} diubah di
          tengah prosedur, atau jalur keluar prosedur melompati
          \lstasm|POP BP|, sehingga epilogue memulihkan nilai yang
          salah.
\end{enumerate}
